% problem_id: S001-lemma-curves-diagonal-separated-fp
% source_id: S001 (Stacks Project)
% source_file: moduli-curves.tex
% statement_locator: moduli-curves.tex lines 291--295
% proof_locator: moduli-curves.tex lines 297--416
% env: lemma
% id_note: label 在本来源内唯一
% pairing: tight (verified)
% status: complete (statement + author proof verbatim + reason + locators)
%
\begin{lemma}
\label{lemma-curves-diagonal-separated-fp}
The diagonal of $\Curvesstack$ is separated
and of finite presentation.
\end{lemma}

\begin{proof}
Recall that $\Curvesstack$ is a limit preserving algebraic stack, see
Quot, Lemma \ref{quot-lemma-curves-limits}.
By Limits of Stacks, Lemma \ref{stacks-limits-lemma-limit-preserving-diagonal}
this implies that
$\Delta : \Polarizedstack \to \Polarizedstack \times \Polarizedstack$
is limit preserving. Hence $\Delta$ is locally of finite presentation
by Limits of Stacks, Proposition
\ref{stacks-limits-proposition-characterize-locally-finite-presentation}.

\medskip\noindent
Let us prove that $\Delta$ is separated. To see this, it suffices to show
that given a scheme $U$ and two objects $Y \to U$ and $X \to U$ of
$\Curvesstack$ over $U$, the algebraic space
$$
\mathit{Isom}_U(Y, X)
$$
is separated. This we have seen in
Moduli Stacks, Lemmas \ref{moduli-lemma-Mor-s-lfp} and
\ref{moduli-lemma-Isom-in-Mor} that the target is
a separated algebraic space.

\medskip\noindent
To finish the proof we show that $\Delta$ is quasi-compact. Since
$\Delta$ is representable by algebraic spaces, it suffices to check
the base change of $\Delta$ by a surjective smooth morphism
$U \to \Curvesstack \times \Curvesstack$ is quasi-compact
(see for example Properties of Stacks, Lemma
\ref{stacks-properties-lemma-check-property-covering}).
We choose $U = \coprod U_i$ to be a disjoint union of affine opens
with a surjective smooth morphism
$$
U \longrightarrow
\textit{PolarizedCurves} \times \textit{PolarizedCurves}
$$
Then $U \to \Curvesstack \times \Curvesstack$ will be surjective
and smooth since $\textit{PolarizedCurves} \to \Curvesstack$
is surjective and smooth by Lemma \ref{lemma-polarized-curves-over-curves}.
Since $\textit{PolarizedCurves}$ is limit preserving
(by Artin's Axioms, Lemma \ref{artin-lemma-fibre-product-limit-preserving}
and Quot, Lemmas \ref{quot-lemma-curves-limits},
\ref{quot-lemma-polarized-limits}, and
\ref{quot-lemma-spaces-limits}), we
see that $\textit{PolarizedCurves} \to \Spec(\mathbf{Z})$ is locally of
finite presentation, hence $U_i \to \Spec(\mathbf{Z})$ is
locally of finite presentation
(Limits of Stacks, Proposition
\ref{stacks-limits-proposition-characterize-locally-finite-presentation}
and Morphisms of Stacks, Lemmas
\ref{stacks-morphisms-lemma-composition-finite-presentation} and
\ref{stacks-morphisms-lemma-smooth-locally-finite-presentation}).
In particular, $U_i$ is Noetherian affine. This reduces us to the
case discussed in the next paragraph.

\medskip\noindent
In this paragraph, given a Noetherian affine scheme $U$ and two objects
$(Y, \mathcal{N})$ and $(X, \mathcal{L})$
of $\textit{PolarizedCurves}$ over $U$, we show the algebraic space
$$
\mathit{Isom}_U(Y, X)
$$
is quasi-compact. Since the connected components of $U$ are open and closed
we may replace $U$ by these. Thus we may and do assume $U$ is connected.
Let $u \in U$ be a point. Let $Q$, $P$ be the Hilbert polynomials
of these families, i.e.,
$$
Q(n) = \chi(Y_u, \mathcal{N}_u^{\otimes n})
\quad\text{and}\quad
P(n) = \chi(X_u, \mathcal{L}_u^{\otimes n})
$$
see Varieties, Lemma \ref{varieties-lemma-numerical-polynomial-from-euler}.
Since $U$ is connected and since
the functions
$u \mapsto \chi(Y_u, \mathcal{N}_u^{\otimes n})$ and
$u \mapsto \chi(X_u, \mathcal{L}_u^{\otimes n})$
are locally constant (see 
Derived Categories of Schemes, Lemma
\ref{perfect-lemma-chi-locally-constant-geometric})
we see that we get the same Hilbert polynomial in every point of $U$.
Set
$$
\mathcal{M} = \text{pr}_1^*\mathcal{N}
\otimes_{\mathcal{O}_{Y \times_U X}} \text{pr}_2^*\mathcal{L}
$$
on $Y \times_U X$. Given $(f, \varphi) \in \mathit{Isom}_U(Y, X)(T)$
for some scheme $T$ over $U$ then for every $t \in T$ we have
\begin{align*}
\chi(Y_t, (\text{id} \times f)^*\mathcal{M}^{\otimes n})
& =
\chi(Y_t,
\mathcal{N}_t^{\otimes n} \otimes_{\mathcal{O}_{Y_t}}
f_t^*\mathcal{L}_t^{\otimes n}) \\
& =
n\deg(\mathcal{N}_t) + n\deg(f_t^*\mathcal{L}_t) +
\chi(Y_t, \mathcal{O}_{Y_t}) \\
& =
Q(n) + n\deg(\mathcal{L}_t) \\
& =
Q(n) + P(n) - P(0)
\end{align*}
by Riemann-Roch for proper curves, more precisely by
Varieties, Definition \ref{varieties-definition-degree-invertible-sheaf} and
Lemma \ref{varieties-lemma-degree-tensor-product}
and the fact that $f_t$ is an isomorphism.
Setting $P'(t) = Q(t) + P(t) - P(0)$ we find
$$
\mathit{Isom}_U(Y, X) =
\mathit{Isom}_U(Y, X) \cap \mathit{Mor}^{P', \mathcal{M}}_U(Y, X)
$$
The intersection is an intersection of open subspaces of
$\mathit{Mor}_U(Y, X)$, see
Moduli Stacks, Lemma \ref{moduli-lemma-Isom-in-Mor} and
Remark \ref{moduli-remark-Mor-numerical}.
Now $\mathit{Mor}^{P', \mathcal{M}}_U(Y, X)$
is a Noetherian algebraic space as it is of finite
presentation over $U$ by
Moduli Stacks, Lemma \ref{moduli-lemma-Mor-qc-over-base}.
Thus the intersection is a Noetherian algebraic space too
and the proof is finished.
\end{proof}
